% DELTA23: full proof draft, root, 2026-11-21 CET.
\subsection{Integral covariance of the five-strand symplectic representation}

\paragraph{D23.T1 (integral specialization of the rigidity method).}
Let $J$ be the alternating $4\times4$ integral matrix with
$J_{i,i+1}=1$, $J_{i+1,i}=-1$ and all other entries zero.
Put $t_v=I+vv^{\mathsf t}J$ and $t_i=t_{e_i}$.
The assignment $q(\sigma_i)=t_i$ defines a continuous homomorphism
$q:\widehat B_5\to\operatorname{Sp}(J,\widehat{\mathbf Z})$.
For every multiplier-one element $h\in\widehat{GT}$ there is a unique
$V_h\in\operatorname{Sp}(J,\widehat{\mathbf Z})$ such that
\[
V_he_1=e_1,\qquad q\Phi_h=\operatorname{Ad}(V_h)q.
\]
The map $h\mapsto V_h$ is a continuous homomorphism.
This statement concerns the genuine five-strand braid action; it does
not assert that this representation descends to the five-punctured
sphere group.

\begin{proof}
The determinant of $J$ is one. The transvections preserve $J$.
For $c=v^{\mathsf t}Jw$, direct multiplication gives
\[
t_vt_w-t_wt_v=c(vw^{\mathsf t}+wv^{\mathsf t})J,
\qquad
t_vt_wt_v-t_wt_vt_w=(1-c^2)(vv^{\mathsf t}-ww^{\mathsf t})J.
\tag{D23.T1.1}
\]
The entries of $J$ therefore give the braid and distant commutation
relations for the $t_i$. The discrete homomorphism extends to the
profinite completion, since reduction modulo each integer has finite
target.

Fix a prime $\ell$. The standard GT action has
$\Phi_h(\sigma_1)=\sigma_1$ and sends every other $\sigma_i$ to a
specified conjugate of itself. Hence
$q_\ell\Phi_h(\sigma_i)=t_{v_i}$ with $v_i$ primitive over
$\mathbf Z_\ell$ and $v_1=e_1$. The $v_i$ can be chosen continuously
using the specified conjugators in the GT braid formula.

We may first work over the field $\mathbf Q_\ell$. If $t_v,t_w$
commute, the first identity in (D23.T1.1) implies $c=0$.
Indeed, if $c\ne0$, it would give $vw^{\mathsf t}+wv^{\mathsf t}=0$.
Choose a linear functional $L$ with $L(v)=1$ and multiply this matrix
on the left by $L$; then $w=-L(w)v$, which implies $c=0$, a
contradiction. For a pair satisfying the braid relation, the second
identity shows that either $c^2=1$ or $t_v=t_w$.

An equality $t_{v_i}=t_{v_{i+1}}$ is impossible. A neighboring
generator $t_{v_{i+2}}$, when present, commutes with $t_{v_i}$ and
braids with $t_{v_{i+1}}$. Two commuting invertible elements which
braid are equal: $a^2b=b^2a$ implies $a=b$. Equality thus propagates
to the right, and the identical argument propagates it to the left.
All four generators would have equal image. But $\Phi_h$ is an
automorphism of $\widehat B_5$, so $q_\ell\Phi_h$ has exactly the
same image as $q_\ell$, and the latter is nonabelian, as witnessed by
$t_1,t_2$. This is a contradiction.

Consequently $c_i=v_i^{\mathsf t}Jv_{i+1}$ equals $1$ or $-1$,
and the pairings of distant vectors are zero. Define
$\epsilon_1=1$ and $\epsilon_{i+1}=\epsilon_i c_i$ and replace
$v_i$ by $w_i=\epsilon_i v_i$. The transvections do not change and
$w_i^{\mathsf t}Jw_{i+1}=1$. Thus the column matrix
$V_\ell=(w_1,w_2,w_3,w_4)$ satisfies
$V_\ell^{\mathsf t}JV_\ell=J$. Taking determinants shows
$\det(V_\ell)^2=1$, so $V_\ell$ is invertible over
$\mathbf Z_\ell$. It fixes $e_1$ and conjugates every $t_i$ to
$t_{v_i}$.

For uniqueness, the simultaneous centralizer of the $t_i$ in
$\operatorname{GL}_4(\mathbf Q_\ell)$ is the scalar group. To see
this, commuting with $t_i-I$ preserves its image $\mathbf Q_\ell e_i$,
so a centralizing matrix is diagonal. Its commutation with the
nonzero $(i,i+1)$ entries of $t_i-I$ equates all adjacent diagonal
entries. A scalar symplectic matrix is $\pm I$, and fixing $e_1$
chooses $I$.

The product $V_h=(V_\ell)_\ell$ lies in the indicated adelic
symplectic group. The displayed formulas using the original
continuous $v_i$ and their pairings prove its continuity, prime by
prime, including at $2$ and $5$. Both $V_{hk}$ and $V_hV_k$ fix
$e_1$ and intertwine $q\Phi_h\Phi_k$ with $q$. Uniqueness gives
$V_{hk}=V_hV_k$.
\end{proof}

\paragraph{D23.T2 (the exact integral order-five centralizer).}
Put $C=t_1t_2t_3t_4$. Then
\[
C=\begin{pmatrix}0&0&0&1\\-1&0&0&1\\0&-1&0&1\\0&0&-1&1\end{pmatrix},
\quad C^5=-I,\quad \det(XI-C)=X^4-X^3+X^2-X+1.
\]
For every commutative ring $R$, the endomorphism centralizer of
$C$ in $M_4(R)$ is $R[C]\simeq R[\zeta_5]$, with
$\zeta_5$ represented by $-C$. Its symplectic centralizer is
\[
\operatorname{Cent}_{\operatorname{Sp}(J,R)}(C)
=\{z\in R[\zeta_5]:z\overline z=1\},
\qquad \overline{\zeta_5}=\zeta_5^{-1}.
\]
\begin{proof}
Multiplication gives the matrix and polynomial identities. The four
vectors $e_1,Ce_1,C^2e_1,C^3e_1$ are respectively
$e_1,-e_2,e_3,-e_4$, so they form a basis over every $R$.
Any endomorphism commuting with $C$ is determined by its value on
$e_1$; writing that value in this basis shows that the endomorphism
is a unique polynomial of degree at most three in $C$. This proves
the asserted algebra centralizer without any restriction on the
residue characteristic.
For the symplectic adjoint $X^*=J^{-1}X^{\mathsf t}J$ one has
$C^*=C^{-1}$, since $C$ is symplectic. On the commutative algebra
$R[C]$ the adjoint is therefore exactly the displayed involution.
The symplectic condition is $z^*z=1$. It already implies that $z$
is a unit, and is precisely the norm-one equation.
\end{proof}
